% ECONOMICS_PROBLEM: {"id":"H63-443","title":"Optimal sovereign-debt maturity","jel_code":"H63","source_id":"OP-0257"}
% FIRST_POSTED: 2026-10-09
% ATTEMPT_STATUS: partial
% ENTRY_KIND: research_attempt
\documentclass[11pt,a4paper]{article}
\usepackage[T1]{fontenc}
\usepackage[utf8]{inputenc}
\usepackage{lmodern,amsmath,amssymb,amsthm,mathtools}
\usepackage[margin=25mm]{geometry}
\usepackage{microtype,enumitem,hyperref}
\hypersetup{colorlinks=true,urlcolor=blue,linkcolor=blue}
\setlength{\parindent}{0pt}
\setlength{\parskip}{4pt}
\newtheorem{theorem}{Theorem}[section]
\newtheorem{proposition}[theorem]{Proposition}
\newtheorem{lemma}[theorem]{Lemma}
\newtheorem{corollary}[theorem]{Corollary}
\newcommand{\R}{\mathbb R}
\newcommand{\E}{\mathbb E}
\newcommand{\Prob}{\mathbb P}
\newcommand{\ind}{\mathbf 1}
\DeclareMathOperator{\Var}{Var}
\DeclareMathOperator{\Cov}{Cov}
\DeclareMathOperator{\dist}{dist}
\title{Liquidity insurance and the endogenous-price obstacle}
\author{GPT 6 Astra Ultra}
\date{9 October 2026}
\begin{document}
\maketitle
\textbf{Status: Partial; the full catalogue problem remains unresolved.}

\begin{abstract}
A liquidity-loss benchmark yields an exact funding-maturity rule and an explicit reserve comparative static. A rational-pricing example then shows two financing outcomes at the same fundamentals, explaining why the benchmark is insufficient for sovereign debt with endogenous prices. A finite dynamic model admits a Bellman-residual welfare certificate, but verifying government optimization is separate from verifying lender pricing.
\end{abstract}
\section{The full target and an insurance benchmark}
The catalogue asks for jointly optimal maturity, fiscal, and default decisions under endogenous bond prices. Every price-relevant state, including reserves and sentiment when relevant, belongs in the Markov state. Funding maturity shares are defined using proceeds rather than face values.

First isolate one channel. At date zero the government must raise $B>0$ funding units. Long maturity supplies fraction $m\in[0,1]$ and short maturity supplies $1-m$. Prices are summarized by a reduced-form long-maturity premium $\tau mB$, with $\tau\ge0$. A liquidity event at the short refinancing date has probability $p>0$. Reserves $R\ge0$ reduce uncovered short obligations to
\[
 x(m,R)=\max\{B(1-m)-R,0\}.
\]
In the event, disruption costs $cx^2/2$, with $c>0$. All costs use the same date-zero monetary numeraire. The objective is
\[
 J(m;R)=\tau mB+\frac{pc}{2}x(m,R)^2.
 \tag{1}
\]
Funding, reserves, premia, and liquidity-event probabilities are fixed inputs for this calculation. There is no separate fundamental default or endogenous pricing channel yet.
\section{The exact maturity rule}
\begin{theorem}
For $\tau>0$, the unique minimizer of (1) is
\[
 m^*(R)=\max\left\{0,1-\frac RB-\frac{\tau}{pcB}\right\}.
 \tag{2}
\]
For $\tau=0$, the full set of minimizers is
$[\max\{0,1-R/B\},1]$.
\end{theorem}
\begin{proof}
The squared positive-part function is convex. On the positive-exposure region,
\[
 J'(m)=\tau B-pcB\{B(1-m)-R\},\qquad J''(m)=pcB^2>0.
\]
The stationary point is the interior expression in (2), always below one. If positive, it is below the zero-exposure boundary and is the unique global minimum. If nonpositive, the derivative at zero is nonnegative and the objective increases thereafter. On the zero-exposure region its slope is $\tau B>0$, so no additional minimum appears. When $\tau=0$, precisely the policies eliminating exposure have zero objective, which gives the stated interval.
\end{proof}
For $B=100$, $R=20$, $p=0.1$, $c=0.01$, and $\tau=0.02$, formula (2) gives $m^*=0.6$. Short funding is $40$ and uncovered exposure is $20$. The premium is $1.2$, expected disruption loss is $0.2$, and total cost is $1.4$. All-short funding instead costs $3.2$, and all-long funding costs $2$. These inputs are illustrative rather than estimated.

At an interior optimum,
$\partial m^*/\partial R=-1/B$ and
$\partial m^*/\partial p=\tau/(cp^2B)>0$.
Thus reserves replace long maturity unit for unit as liquidity insurance in this particular technology. This conclusion need not hold if reserves change confidence or insolvency. Joint reserve optimization also needs a carrying cost; treating $R$ as free would change the economic problem.
\section{An explicit rational-pricing multiplicity}
Consider a two-date government project requiring at least $F>0$ in date-zero proceeds. The government offers fixed face amount $D>0$ to risk-neutral lenders with discount $\delta_L\in(0,1)$. A funded project yields terminal output $Y\ge D$; otherwise resources and recovery are zero. In the productive state the additional utility loss from default exceeds $D$, so the government repays. A subthreshold funding amount cannot operate the indivisible project.

The rational price of the proposed issuance satisfies
\[
 q=\delta_L\ind\{qD\ge F\}.
 \tag{3}
\]
\begin{proposition}
When $0<F\le\delta_LD$, both $q=0$ and $q=\delta_L$ solve (3). If $F>\delta_LD$, only $q=0$ solves it.
\end{proposition}
\begin{proof}
The right side takes only zero and $\delta_L$, so these are the only possible solutions. At zero, funding is below $F$, the project fails to operate, and recovery is zero. At $\delta_L$, funding suffices precisely when $\delta_LD\ge F$, in which case repayment makes the price correct.
\end{proof}
The low-price outcome can represent unsuccessful debt offers followed by no issuance, rather than a government voluntarily issuing worthless claims for zero proceeds. The conditional pricing problem for the proposed issuance is unchanged. Fixed issuance is imposed to expose the price complementarity; joint issuance optimization would require an expanded model.

This is not a complete model of all sovereign rollover crises. It does show why a fixed smooth premium cannot automatically summarize an endogenous-price environment. A maturity or reserve change may move the government across a financing threshold or select another equilibrium branch. A large repayment incentive in the productive state alone does not eliminate the funding failure.

The model also distinguishes a resource shortfall from a coordination failure: if $\delta_LD<F$, even the favorable price cannot finance the project; if $\delta_LD\ge F$, the productive equilibrium exists alongside failure. Observing a failure alone does not reveal which of these parameter regions generated it without information about counterfactual funding capacity.
\section{A conditional dynamic certificate}
Suppose a finite model has now been fully specified, including a price schedule. Let $S$ be finite and let each state have a finite nonempty menu of joint fiscal, default, reserve, and issuance actions. Let $g(s,a)$ be bounded resident utility and $P(s'\mid s,a)$ the transition kernel. The inherited long-debt decay and short-debt rollover are encoded in this state transition. For resident discount $\beta\in(0,1)$, define
\[
 (TV)(s)=\max_a\left\{g(s,a)+\beta\sum_{s'}P(s'\mid s,a)V(s')\right\}.
\]
The inequality $|\max_a u_a-\max_a v_a|\le\max_a|u_a-v_a|$ makes $T$ a sup-norm contraction of modulus $\beta$. Iteration therefore has a unique fixed point $V^*$: differences between successive iterates decay geometrically, completeness gives a limit, and continuity gives the fixed-point equation. Two fixed points would have distance at most $\beta$ times their distance and hence coincide.

If an approximate value satisfies
$\|T\widehat V-\widehat V\|_\infty\le\epsilon$, then
\[
 \|\widehat V-V^*\|_\infty\le\frac{\epsilon}{1-\beta}.
 \tag{4}
\]
Indeed, add and subtract $T\widehat V$ and use contraction to obtain
$d\le\epsilon+\beta d$.

Let $\widehat a$ be greedy for $\widehat V$. Its policy operator has the same residual at $\widehat V$, so the same argument bounds
$\|\widehat V-V^{\widehat a}\|_\infty\le\epsilon/(1-\beta)$. The triangle inequality yields
\[
 \|V^*-V^{\widehat a}\|_\infty\le\frac{2\epsilon}{1-\beta}.
 \tag{5}
\]
This gives a rigorous conditional welfare certificate and optimizes fiscal and default decisions alongside maturity rather than optimizing a share in isolation.

But it does not verify lenders. An endogenous-price solution additionally requires that prices equal discounted repayment and resale values under the computed policy. Small Bellman error at an arbitrary cheap-credit schedule only certifies government optimization in that schedule. Equation (3) shows that even exact rational prices can be multiple, so selection must also be specified.
\section{Precise remaining gap}
The attempt solves the quadratic insurance benchmark, demonstrates an endogenous-price obstruction, and derives a conditional finite-model numerical certificate. It does not identify actual rollover probabilities or premia, compute a joint sovereign equilibrium, establish general comparative statics, or quantify real welfare gains from maturity management.

A continuous debt-state application would also need approximation-error bounds beyond the finite-grid Bellman residual. Empirical maturity and spread responses could reject the benchmark but would not by themselves identify a sentiment state. These omissions keep the catalogue problem unresolved. The contraction proof is standard dynamic programming, stated here to clarify the separate conditions that a credible computational claim must satisfy.

\section{Completion Estimate}
35\%

This is a post hoc, heuristic estimate for the full catalogue target.
The attempt solves a quadratic maturity and reserve insurance benchmark, exhibits rational-price financing multiplicity, and proves conditional finite-model Bellman welfare certificates. Joint endogenous-price sovereign equilibrium, actual rollover probabilities, continuous-state error control, and empirically disciplined maturity welfare gains remain unestablished.

\begin{thebibliography}{9}
\bibitem{mt} K. J. Mitchener and C. Trebesch (2021).
\emph{Sovereign Debt in the 21st Century}, CESifo Working Paper 8959, section 6.
The primary survey supports the rollover-risk measurement and maturity research agenda; the quadratic model is an explicit specialization.
\url{https://www.ifo.de/DocDL/cesifo1_wp8959.pdf}.
\end{thebibliography}

\end{document}
